\section{Architectural Decision: Dedicated Slot for STARK-FRI}
\label{sec:architectural-decision}

\subsection{The 2026-06-03 slot relocation}

Pre-rename state: the Plonky3-fork STARK / FRI verifier lived at
slot \texttt{0x012205} under the package name \texttt{p3q}. The
package contained the line \emph{``P3Q strict-PQ STARK verifier''}
and the magic-header tag \texttt{"P3Q1"} was threaded through every
proof artifact (EasyCrypt, Lean, Jasmin, dudect).

Per Hanzo Crypto Suite \S~5.2~\cite{hanzocrypto} and the roadmap
\S~B.11, \textbf{P3Q canonically means ``Post-Quantum Pulsar
Proof''}: the on-chain unified PQ \emph{threshold-signature}
verifier dispatch (Pulsar~/ Corona~/ Magnetar). The STARK / FRI
dispatch occupying slot \texttt{0x012205} under the same name was
an aliasing error: STARK proof verification and PQ signature
verification are two different primitives serving disjoint use
cases.

The 2026-06-03 decomplection makes this explicit by moving each
concern to its own slot:

\begin{table}[h]
\centering
\small
\begin{tabular}{@{}llll@{}}
\toprule
\textbf{Surface} & \textbf{Pre-rename} & \textbf{Post-rename} & \textbf{Reason} \\
\midrule
Source path & \texttt{lux/precompile/p3q/} & \texttt{lux/precompile/starkfri/} & primitive name \\
Slot & \texttt{0x012205} & \texttt{0x012220} & disjoint surface \\
Go package & \texttt{package p3q} & \texttt{package starkfri} & value-typed naming \\
Address symbol & \texttt{ContractP3QVerifyAddress} & \texttt{ContractStarkFRIVerifyAddress} & matches package \\
Singleton type & \texttt{P3QVerifyPrecompile} & \texttt{StarkFRIVerifyPrecompile} & matches package \\
Config key & \texttt{p3qVerify} & \texttt{starkfriVerify} & matches package \\
Wire-bytes magic & \texttt{"P3Q1"} & \texttt{"P3Q1"} (preserved) & proof-artifact stability \\
\bottomrule
\end{tabular}
\caption{Surfaces touched in the 2026-06-03 rename. Slot moves;
package and symbol names follow the package; the wire-bytes magic
header is preserved verbatim because every proof artifact references
it.}
\label{tab:rename-surfaces}
\end{table}

After the rename:

\begin{itemize}[leftmargin=2em,nosep]
\item \texttt{0x012205} is LP-218's unified PQ \emph{signature}
  verifier (P3Q-Pulsar~/ P3Q-Corona~/ P3Q-Magnetar -- kind-byte
  dispatched per LP-220).
\item \texttt{0x012220} is LP-221's STARK-FRI \emph{proof}
  verifier (this LP).
\end{itemize}

The two slots are orthogonal. There is no aliasing, no overlap, no
shared dispatch table. Both can be called from a single ZK-rollup
batch validator and the two costs sum cleanly:
$\sim$\num{50000} gas for Pulsar verify $+$ $\sim$\num{500000}
gas for STARK verify $\approx$ \num{550000} gas per ZK-rollup
batch at typical parameters.

\subsection{Why \texttt{"P3Q1"} survives the rename}

The 4-byte tag \texttt{"P3Q1"} is a \emph{wire-format byte string},
not a name claim. Renaming it would invalidate the byte-identical
proof artifacts (EasyCrypt theories at admit-budget $0\,/\,0$, Lean
theorems, Jasmin source, dudect harness) that reference the header
verbatim. A future v2 wire format MAY rename the magic to
\texttt{"SFR1"} (STARK-FRI Rev~1) once the proof artifacts and the
artifact-budget tracker (\texttt{scripts/checks/ec-admits.sh}) are
refreshed in lockstep. Today, breaking proof-artifact byte equality
for a cosmetic rename has no upside.

This is the \emph{values, not places} discipline: the magic header
is the value at the wire, not a label about which file it lives in;
renaming the file does not rename the wire-bytes value, and the
proof artifacts that reference the wire-bytes value stay byte-stable.

\subsection{Lockstep updates across the codebase}

The rename touched five downstream surfaces in the same atomic
landing. These are recorded for self-containment; the normative
surface remains the precompile package itself.

\begin{itemize}[leftmargin=2em,nosep]
\item \texttt{lux/geth/core/vm/lux\_precompiles.go} -- EVM wiring,
  registry entry updated to
  \texttt{starkfri.StarkFRIVerifyPrecompile} at \texttt{0x012220}.
\item \texttt{lux/chains/evm/main.go} -- chain import shim.
\item \texttt{lux/evm/precompile/registry/registry.go} --
  precompile registry.
\item \texttt{lux/crypto/p3q/ct/dudect/verify\_ct.go} -- dudect
  harness; cross-tree rename to
  \texttt{lux/crypto/starkfri/ct/dudect/} pending (see
  \S~\ref{sec:constant-time}).
\item LP-218~\S``Verifier reference'' -- inline reference to the
  STARK-FRI precompile now points at \texttt{0x012220} rather than
  the old \texttt{0x012205} co-location.
\end{itemize}

\subsection{Why a Plonky3 fork}

The reference backend at \texttt{lux/p3q/crates/p3q-verifier} is a
Plonky3~\cite{plonky3} fork. The fork strips:

\begin{itemize}[leftmargin=2em,nosep]
\item KZG polynomial commitments (require pairings, classical-only).
\item BN254 / BLS12-381 elliptic-curve backends.
\item All non-strict-PQ FRI configurations (e.g.\ small-field
  Reed-Solomon parameter choices that don't meet the soundness
  target of $\le 2^{-100}$ under the chosen blowup factor and query
  count).
\end{itemize}

Plonky3 upstream is dual-tracked: a strict-PQ profile (cSHAKE256~/
Goldilocks~/ FRI only) and a classical profile (KZG~/ Poseidon~/
BN254). LP-221's fork pins the strict-PQ profile and removes the
classical surface entirely. The two reasons:

\begin{enumerate}[leftmargin=2em,nosep]
\item \textbf{Audit surface.} An EVM precompile is a permanent
  reservation in block validation. Every line of code reachable
  from the precompile entry point is in the audit boundary. A
  dual-tracked upstream with both classical and PQ surfaces
  doubles the audit boundary without doubling the security
  guarantee -- only the strict-PQ surface is reachable in
  production.
\item \textbf{Misuse resistance.} A future caller cannot
  accidentally instantiate the classical surface by passing the
  wrong configuration bytes; the classical surface does not exist
  in the binary. The wire-format version byte
  (\texttt{VersionV1=0x01}) selects \emph{which strict-PQ profile}
  the verifier expects; it does not select between classical and
  PQ.
\end{enumerate}

Tracking the Plonky3 upstream beyond this point is a
\emph{rebase-and-re-strip} discipline: each upstream release the
fork rebases, re-applies the strip patches, re-runs the EasyCrypt
admit budget check, and re-tests against the proof corpus. The fork
maintainer is one person on the cryptography team; the rebase is
not a routine release-engineering concern.

\subsection{What this is not}

It is worth being explicit about two design decisions LP-221 does
\emph{not} make:

\begin{enumerate}[leftmargin=2em,nosep]
\item \textbf{It is not a kind-byte-dispatched family.} Unlike
  LP-218~/ LP-220 at slot \texttt{0x012205} where a single precompile
  dispatches over $\{0x01, 0x02, 0x03\}$ to three signature
  verifiers, LP-221 is a single-primitive precompile. The
  \texttt{VersionV1=0x01} version byte exists so a future LP can
  graft a v2 wire format (post-Plonky3-upstream-rotation or a
  recursive variant) without re-allocating the slot, but the
  version is not a primitive-family discriminator. The slot is
  reserved for STARK-FRI; if a fundamentally different SNARK family
  enters the Lux stack (a recursive HyperPlonk variant, say) it
  lands at its own slot inside the reserved range
  \texttt{0x012220..0x01222F}.
\item \textbf{It is not a backwards-compatible upgrade of a prior
  STARK precompile.} There is no pre-LP-221 STARK verifier
  deployment in the wild. LP-0120's historical row at
  \texttt{0x012205} for ``Plonky3 PQ-only STARK'' was a
  pre-genesis placeholder that never reached any network. The
  activation marker \texttt{2025-12-25T16:20:00-08:00} places the
  LP-221 ABI at network genesis of the new final Lux network. The
  pre-Quasar Edition Lux network (2020--2025) had no STARK-FRI
  precompile; there is nothing to migrate, no shim, no replay
  window.
\end{enumerate}

\subsection{Activation marker}

\begin{verbatim}
activates:      2025-12-25T16:20:00-08:00
activates-unix: 1766708400
\end{verbatim}

The STARK-FRI precompile at slot \texttt{0x012220} is callable from
the genesis block of the new final Lux network.
